\newpage
\section*{Summary} % Use \section* instead of \chapter*
\addcontentsline{toc}{chapter}{Summary} % Add to ToC at chapter level
\markboth{Summary}{} % Set header text

\vspace{0.3cm} % Adjust space after title

\subsection*{English version}
Antimicrobial resistance has been declared by the World Health Organization as one of the top global health concerns of our current century. It has been associated with more than one million deaths in 2019 and it is projected to cause up to 10 million deaths per year by 2050 if action is not taken. Antimicrobial resistance is closely linked to the misuse and overuse of antibiotics, which trigger natural evolutionary processes in which bacteria eventually develop mechanisms to evade these compounds. As a consequence, antibiotics progressively lose their effectiveness, leaving humanity increasingly vulnerable against infections. Most conventional antibiotics act by binding to specific cellular targets and inhibiting microbial proliferation. Since developing new antibiotics that display a similar mode of action will eventually make bacteria to develop resistance against them, a change in the paradigm is needed.

In this regard, bacterial membranes emerge as a promising therapeutic target, since the lipid compositions of bacterial and mammalian cells differ significantly, providing opportunities for selective targeting strategies. More specifically, eukaryotic membranes present a content of phosphatidylcholine lipids close to the 50\%, followed by the anionic phosphatidylserine lipids, which are especially abundant in the inner leaflet. This asymmetric lipid distribution in eukaryotic membranes results in an outer leaflet that is overall neutrally charged. In contrast, bacterial membranes are mainly composed of the anionic phosphatidylglycerol and the neutral phosphatidylethanolamine, which together confer a net negative charge to the membrane surface. Taking advantage of these differences, molecules that preferentially interact with negatively charged membranes represent attractive candidates for selective antibacterial activity.

Antimicrobial peptides, in particular membrane-targeted antimicrobial peptides, have therefore attracted considerable attention, as they are hypothesized to selectively target negatively charged membranes. As part of the innate immune system in many living beings, they exhibit a wide structural diversity and are associated with multiple biological functions. Among these, antimicrobial peptides display activity against a broad range of microorganisms, including Gram-positive and Gram-negative bacteria. Their biological activity is commonly associated to their amphipathic nature: the cationic regions promote initial electrostatic attraction to negatively charged membranes, while hydrophobic residues facilitate insertion into the non-polar core of the lipid bilayer. These combined driving forces lead to peptide accumulation, facilitating insertion, disturbing the membrane integrity, and ultimately leading to membranolytic effects that trigger cell death.

Nevertheless, the use of antimicrobial peptides as therapeutic agents is frequently hindered by multiple factors, including poor metabolic stability, limited bioavailability, and potential toxicity toward host cells. These limitations significantly restrict their use as conventional drugs despite their potent antimicrobial activity. Peptide cyclization has therefore emerged as a promising strategy to overcome these limitations, as cyclic versions of antimicrobial peptides have already demonstrated enhanced stability and improved biological performance.
%%%%%% CONTINUE
Particularly interesting in this context are \textit{D,L}-$\alpha$-Cyclic Peptides, which have been shown to be largely inactive as isolated monomeric units, while exhibiting biological activity upon self-assembly into nanotubular structures. When formed by an even number of alternating \textit{D}- and \textit{L}-$\alpha$-amino acids, they adopt a flat, ring-shaped structure with the amino acid side-chains projected outward. In this unique geometry the amide groups of the backbone are oriented perpendicular to the plane of the ring, allowing the formation of hydrogen bonds with neighboring cyclic peptides. The resulting self-assembled cyclic peptide nanotubes have been shown to exhibit different behaviors in the presence of lipid bilayers, depending on their amino acid composition. Cyclic peptides bearing hydrophobic side-chains can form nanotubes that span the thickness of lipid bilayers, leading to transmembrane channels capable of ion transport. In contrast, amphiphatic cyclic peptides tend to adsorb onto the membrane surface, where they can exert antimicrobial activity. Despite these promising properties, the molecular mechanisms governing cyclic peptide self-assembly and membrane interactions are not fully understood. Experimental investigation of these processes often require techniques with atomic-level spatial and temporal resolution, which are not accessible to most experimental approaches. In this context, Molecular Dynamics simulations provide a powerful computational framework to investigate these processes by offering a time-resolved, molecular-level description of peptide structures, dynamics, and interactions.

However, Molecular Dynamics simulations of systems comprising lipid membranes and multiple cyclic peptides are severely limited by computational cost, particularly when studying collective phenomena such as self-assembly, membrane insertion, or aggregation. To overcome these limitations, Coarse-Grained Molecular Dynamics approaches offer a practical alternative by reducing the number of degrees of freedom and enabling simulations over longer timescales and larger system sizes. In the case of \textit{D,L}-$\alpha$-cyclic peptides, the use of coarse-grained models introduces a critical challenge: their self-assembly is governed by backbone orientation and hydrogen-bond directionality, both of which are intrinsically linked to chirality. Standard coarse-grained force fields, such as Martini, do not explicitly encode this information, leading to unrealistic stacking geometries and inaccurate descriptions of supramolecular assembly. Accurately accounting for chirality at the coarse-grained level becomes therefore essential to faithfully model these systems .At the same time, the adoption of coarse-grained representations enables systematic and high-throughput simulation strategies, allowing the exploration of large numbers of cyclic peptide sequences and membrane interaction scenarios under fully controlled and reproducible conditions. Such studies generate large and homogeneous datasets that describe the structure, dynamics, and membrane interactions of peptides. The availability of these datasets is particularly relevant in the context of Machine Learning. Current predictive models for cyclic peptide activity predominantly rely on sequence-based representations, directly mapping amino acid sequences to biological outcomes. While successful in many cases, these approaches often lack explicit information about the underlying molecular mechanisms governing peptide–membrane interactions. In contrast, Molecular Dynamics simulations provide access to rich, physically meaningful descriptors---such as interaction energies, contact patterns, structural rearrangements, and dynamical properties---that are directly linked to function.

This doctoral thesis is grounded in this synergy between molecular simulation and artificial intelligence. The ultimate goal is to provide new computational methods and tools that support the rational design of cyclic peptides, particularly for antimicrobial applications, through physically informed simulations and data-driven prediction. To achieve this, the work is structured around a series of core objectives:
\begin{enumerate}
    \item The development of a chirality-aware coarse-grained force field that improves the representation of cyclic peptides while enabling microsecond-scale simulations.
    \item The creation of open-source computational tools for the automated generation of three dimensional structures and topologies of cyclic peptide sequences and self-assembled cyclic peptide nanotubes, with the aim of lowering technical barriers and democratizing access to Molecular Dynamics simulations of these systems.
    \item The application of chirality-aware coarse-grained models and computational tools to perform systematic and high-throughput Molecular Dynamics simulations to investigate cyclic peptide self-assembly and their interactions with membrane models.
    \item The extraction, analysis, and organization of simulation-derived descriptors into a public and reusable Molecular Dynamics database, designed to support data-driven analysis and model development.
    \item The development and application of Machine Learning models trained on both experimental data and Molecular Dynamics–derived descriptors to predict cyclic peptide membrane permeability and guide mechanism-informed design strategies.
\end{enumerate}
%%%%%%%%%
The first major achievement of this thesis was the development of MA(R/S)TINI, a novel chirality-aware adaptation of the Martini 2.2 coarse-grained force field, specifically designed to model the self-assembly of \textit{D,L}-$\alpha$-cyclic peptides into self-assembled cyclic peptide nanotubes. This model overcomes several fundamental limitations of the standard Martini force field: the overestimation of inter-peptide distances within nanotubes, its inability to distinguish between parallel and antiparallel $\beta$-sheet arrangements, and the formation of stacked assemblies with unrealistic rotational angles that would result in donor–acceptor mismatches at the atomistic level. The latter limitation is particularly critical, as it leads to an artificial overestimation of cyclic peptide assembly. To evaluate MA(R/S)TINI's performance, three representative cyclic peptide sequences were selected: CP1 (\underline{R}R\underline{K}W\underline{L}W\underline{L}W, amphipathic with proven antimicrobial activity), CP2 (\underline{L}L\underline{L}L\underline{L}L\underline{L}L fully hydrophobic), and CP3 (\underline{L}Q\underline{L}W\underline{L}W\underline{L}W, hydrophobic with ability to form transmembrane nanotubes). These peptides were simulated in the presence of three different membrane models: one mammalian and two bacteria-like membranes. The results were systematically compared with those obtained for the same systems using the standard Martini 2.2 force field.

Large-scale simulations performed at varying cyclic peptide concentrations (20, 40, 60, 80 and 160 units) demonstrated that MA(R/S)TINI produced self-assembly patterns similar to those observed with Martini 2.2, while providing a more physically consistent description of stacking geometries. Importantly, the force field captured the expected interactions of the different cyclic peptides sequences with specific membrane compositions. For example, CP1 formed nanotubes preferentially in the aqueous phase when simulated with the mammalian model, and at the surface of the lipid bilayer when exposed to a bacterial model. In contrast to Martini 2.2, MA(R/S)TINI correctly differentiated between parallel and antiparallel arrangements and restricted rotational angles to physically meaningful values. Metadynamics simulations of cyclic peptide dimers further confirmed that the introduction of additional backbone interaction sites effectively addressed the geometric limitations of the original force field. However, Potential of Mean Force calculations showed an interaction energy almost twice as high in MA(R/S)TINI, which was probably leading to an overestimation of the self-assembly.

The release of a new Martini version, Martini 3, which introduced a revised set of interaction beads, enabled the refinement of this approach through the development of MA(R/S)TINI 3. This new version was designed to retain the geometric advantages of the chirality-aware model while restoring interaction energies closer to those of the parent force field. To ensure a fair comparison between MA(R/S)TINI 3 and its predecessor, the same simulation protocols and test systems were employed. Unbiased Molecular Dynamics demonstrated that MA(R/S)TINI 3 preserved sequence-dependent self-assembly and improved stacking geometries, while producing more dynamic supramolecular assemblies consistent with a reduction in inter-peptide interaction strength. Metadynamics simulations confirmed the improvement upon the original version by reducing the exaggerated interaction energy between cyclic peptides to match that of the original force field. By preserving the strengths of MA(R/S)TINI and correcting its weaknesses, MA(R/S)TINI 3 offers a robust platform for studying cyclic peptide-based supramolecular nanotubes and their interactions with membranes.

In order to facilitate the access to cyclic peptide Molecular Dynamics simulations and three dimensional structure visualization, both original and adapted coarse-grained force fields along with multiple atomistic models were integrated into a tool as part of a publicly available web server under the name CYCLOPEp. The CYCLOPEpBuilder tool provides a user-friendly interface for the creation of cyclic peptide sequences or self-assembled cyclic peptide nanotubes. This tool provides full control over the structure, including: the selection of sequence length ranging from 4 to 12 amino acids, customization of the relative orientation between cyclic peptide pairs, as well as rotation of one sequence relative to its neighbor. After minimal calculations, the platform provides a set of files containing the three dimensional structures of the cyclic peptide(s) and the assembled nanotube, along with the simulation files necessary to set up a Molecular Dynamics simulations. By automating the setup of complex structures and minimizing technical barriers, CYCLOPEpBuilder democratizes access to simulation methods, promoting their adoption within cyclic peptide research.

Building on these tools, a large-scale high-throughput MD study was conducted, resulting in CYCLOPEpDB, a database of coarse-grained Molecular Dynamics simulations of cyclic peptide-membrane systems using MA(R/S)TINI 3. The database includes results for 3855 unique \textit{D,L}-$\alpha$-cyclic octapeptides interacting with a bacterial membrane model. Each sequence was used to build an eight-membered nanotube which was simulated in three orientations relative to the membrane surface. A uniform setup and simulation pipeline ensured consistent conditions across experiments. From each 10-microsecond trajectory, 1006 quantitative descriptors were extracted, covering interaction energies, number of contacts, or distances among others.

As a proof of concept, a subset of six features capturing key aspects of the interaction between cyclic peptides and the membrane was selected from the full feature space. Next, Principal Component Analysis was applied to reduce the dimensionality of the data to two principal components, which were then used for clustering with the OPTICS algorithm. Analysis of the distribution of the six initial features across the three identified clusters revealed distinct classes of membrane interaction, which were categorized as Lipophilic, Intermediate and Lipophobic. Further analysis of these categories revealed correlations with amino acid composition and net charge: Lipophilic cyclic peptides were enriched in basic residues but preferentially had net charges close to zero, Intermediate sequences were mostly highly positively charged with a balanced amino acid profile, whereas Lipophobic cyclic peptides were rich in acidic residues and predominantly carried a net negative charge.

To enable Machine Learning predictions based on the simulation data, each sequence was converted into a graph, based on its coarse-grained topology and force field parameters. Graph nodes encoded bead-level features (mass, charge, coordinates, and Lennard-Jones interaction values), while edges represented bond connectivity and bond lengths. A Message Passing Neural Network trained on this graph representation achieved outstanding performance in both classification task with an accuracy close to 90\% and in regression task predicting a simulation based pseudopermeability score with \(R^2 \approx\)0.92 and a mean squared error of \(0.51\). The used architecture in combination with the MA(R/S)TINI 3 based graphs significantly outperformed baseline models demonstrating the power of combining Molecular Dynamics simulations with graph-based Machine Learning. Moreover, all simulation data and extracted descriptors were made publicly available through integration of CYCLOPEpDB into CYCLOPEp platform, enabling other researchers to made their own analyses and train predictive models.

In parallel, the release of CycPeptMPDB, a database containing experimental membrane permeability data for cyclic peptides, allowed the development of complementary Machine Learning models trained directly on experimental measurements. The aim of this study was to optimize both classification and regression performance by leveraging chemical intuition in the choice of molecular representations, data augmentation strategies, and neural network architectures. First, three types of input representations were systematically evaluated:
\begin{enumerate}
    \item Sequential properties encoding amino acid order with residue-level features.
    \item Cyclic peptide-level physicochemical descriptors.
    \item SMILES-based encoding.
\end{enumerate}
The best performing model, which combined sequential properties and cyclic peptide-level features, was trained on two different data augmentation strategies. The first method exploited the circularity of the sequences by applying cyclic permutations, that is, shifting the ordering of the amino acids within the sequence by one position, adding to the dataset all the possible combinations that represented each sequence. The second method aimed to greatly augment the number of sequences by means of amino acid substitutions. Since the replacement of one amino acid by another may have a huge impact on activity, especially on short peptides, different substitution rules were tested. Some datasets were designed to be more conservative, looking for an improvement in the prediction performance, while others introduced disruptive changes to confuse the model and study the role of proper amino acid mutation. While cyclic permutations showed a slight enhancement of the predictive accuracy, all the studied mutation datasets reduced the performance of the original model. It was observed that, more conservative mutations yielded more accurate models, emphasizing the importance of careful selection of amino acids.

Finally, two architectural modifications were explored to improve the ability of the model to capture the cyclic nature of the sequences. In a first approach, the conventional convolutional layers were replaced by cyclic convolutional layers, specifically designed to account for ring-closure patterns. This modification yielded very similar results to those of the original model, suggesting that standard convolutional layers may already capture the cyclic information without additional modifications. In a second approach, each cyclic permutation of a sequence was processed independently using a convolutional neural network followed by a Long Short-Term Memory module, and the output tensors were aggregated to produce final predictions. Despite a slight increase in computational cost, the model achieved predictive performances comparable to those of more complex deep neural networks published in the literature, suggesting that data quality, rather than model complexity, may be the limiting factor.

To facilitate broad accessibility and practical usage, the best-performing experimental permeability model was deployed as CYCLOPS, a web-based tool integrated into the CYCLOPEp platform. CYCLOPS enables users to predict membrane permeability using the cyclic peptide sequence as input. The sequence builder supports more than 300 amino acids and allows the construction of cyclic peptides ranging from 2 to 15 residues with both head-to-tail and head-to-side-chain cyclization schemes.

In summary, this work highlights the importance of chemically informed molecular representations and cautious augmentation strategies in order to build accurate models for cyclic peptide permeability prediction. Moreover, it provides a user-friendly platform offering accessibility to Machine Learning tools to researchers without expertise.

In conclusion, this doctoral thesis provides CYCLOPEp, an open-source online platform integrating three complementary components aimed at accelerating cyclic peptide research. CYCLOPEpBuilder enables the automated generation of three-dimensional structures and topologies of \textit{D,L}-$\alpha$-Cyclic Peptides and self-assembled cyclic peptide nanotubes, either for simulation or structure visualization. Different levels of resolution are available, including the MA(R/S)TINI coarse-grained models specifically tailored to reproduce cyclic peptide characteristics at a reduced computational cost. CYCLOPEpDB provides an extensive database of coarse-grained Molecular Dynamics simulations of \textit{D,L}-$\alpha$-cyclic octapeptides in the presence of a bacterial membrane model, at the MA(R/S)TINI 3 resolution level. The database includes precomputed analyses and simulation-derived descriptors that can be directly used for data analysis and Machine Learning. Finally, CYCLOPS integrates Machine Learning models trained on experimental cyclic peptide permeability data to predict membrane permeability. Its sequence builder supports more than 300 monomers, enabling the construction of cyclic peptides of varying lengths and cyclization schemes. The results of this thesis offer a solid platform for the study of cyclic peptides using computational tools. We expect CYCLOPEp to serve as a valuable resource for the scientific community by facilitating the access to computational tools and data, and by supporting the integration of computational modeling techniques with wet-lab experiments toward the rational design of cyclic peptides with antimicrobial activity.














\subsection*{Resumo}

A resistencia antimicrobiana foi declarada pola Organización Mundial da Saúde como unha das principais preocupacións de saúde global do noso século. Asociouse con máis dun millón de mortes en 2019 e proxéctase que causará ata 10 millóns de mortes ao ano para 2050 se non se toman medidas. A resistencia antimicrobiana está estreitamente ligada ao mal uso e abuso dos antibióticos, o que desencadea procesos evolutivos naturais nos que as bacterias acaban desenvolvendo mecanismos para evadir estes compostos. Como consecuencia, os antibióticos perden progresivamente a súa eficacia, deixando á humanidade cada vez máis vulnerable fronte ás infeccións. A maioría dos antibióticos convencionais actúan uníndose a dianas celulares específicas e inhibindo a proliferación microbiana. Dado que o desenvolvemento de novos antibióticos que presenten un modo de acción similar fará que as bacterias acaben desenvolvendo resistencia contra eles, é necesario un cambio de paradigma.

Neste sentido, as membranas bacterianas xorden como unha prometedora diana terapéutica, xa que as composicións lipídicas das células bacterianas e de mamíferos difiren significativamente, o que ofrece oportunidades para estratexias de actuación selectiva. Máis especificamente, as membranas eucariotas presentan un contido de lípidos fosfatidilcolina próximo ao 50\%, seguido polos lípidos aniónicos fosfatidilserina, que son especialmente abundantes na capa interna. Esta distribución lipídica asimétrica nas membranas eucariotas dá lugar a unha cara externa que ten, en xeral, carga neutra. Pola contra, as membranas bacterianas están compostas principalmente polo fosfatidilglicerol aniónico e a fosfatidiletanolamina neutra, que en conxunto confiren unha carga neta negativa á superficie da membrana. Aproveitando estas diferenzas, as moléculas que interaccionan preferentemente con membranas cargadas negativamente representan candidatos atractivos para unha actividade antibacteriana selectiva.

Os péptidos antimicrobianos, en particular os péptidos antimicrobianos dirixidos á membrana, atraeron por tanto unha atención considerable, xa que se postula que teñen como diana selectiva as membranas cargadas negativamente. Como parte do sistema inmunitario innato de moitos seres vivos, exhiben unha ampla diversidade estrutural e asócianse con múltiples funcións biolóxicas. Entre estas, os péptidos antimicrobianos mostran actividade contra unha ampla gama de microorganismos, incluíndo bacterias Gram-positivas e Gram-negativas. A súa actividade biolóxica asóciase comunmente á súa natureza anfipática: as rexións catiónicas promoven a atracción electrostática inicial cara ás membranas cargadas negativamente, mentres que os residuos hidrófobos facilitan a inserción no núcleo apolar da bicapa lipídica. Estas forzas impulsoras, combinadas, levan á acumulación do péptido, facilitando a inserción, alterando a integridade da membrana e, finalmente, conducindo a efectos membranolíticos que desencadean a morte celular.

Non obstante, o uso dos péptidos antimicrobianos como axentes terapéuticos vese frecuentemente obstaculizado por múltiples factores, incluíndo unha baixa estabilidade metabólica, unha biodispoñibilidade limitada e unha potencial toxicidade cara ás células hóspede. Estas limitacións restrinxen significativamente o seu uso como fármacos convencionais a pesar da súa potente actividade antimicrobiana. Polo tanto, a ciclación peptídica xurdiu como unha estratexia prometedora para superar estas limitacións, xa que as versións cíclicas dos péptidos antimicrobianos xa demostraron unha estabilidade aumentada e un mellor rendemento biolóxico.

Particularmente interesantes neste contexto son os \textit{D,L}-$\alpha$-Péptidos Cíclicos, que demostraron ser en gran medida inactivos como unidades monoméricas illadas, mentres que exhiben actividade biolóxica tras a súa autoensamblaxe en estruturas nanotubulares. Cando están formados por un número par de \textit{D}- e \textit{L}-$\alpha$-aminoácidos alternos, adoptan unha estrutura plana en forma de anel coas cadeas laterais dos aminoácidos proxectadas cara ao exterior. Nesta xeometría única, os grupos amida do esqueleto oriéntanse perpendicularmente ao plano do anel, permitindo a formación de pontes de hidróxeno con péptidos cíclicos veciños. Os nanotubos de péptidos cíclicos autoensamblados resultantes demostraron exhibir diferentes comportamentos en presenza de bicapas lipídicas, dependendo da súa composición de aminoácidos. Os péptidos cíclicos portadores de cadeas laterais hidrófobas poden formar nanotubos que abarcan o grosor das bicapas lipídicas, dando lugar a canles transmembrana capaces de transportar ións. Pola contra, os péptidos cíclicos anfipáticos tenden a adsorberse na superficie da membrana, onde poden exercer actividade antimicrobiana. A pesar destas propiedades prometedoras, os mecanismos moleculares que gobernan a autoensamblaxe dos péptidos cíclicos e as interaccións coa membrana non se comprenden completamente. A investigación experimental destes procesos require, a miúdo, técnicas con resolución espacial e temporal a nivel atómico, que non son accesibles para a maioría dos enfoques experimentais. Neste contexto, as simulacións de Dinámica Molecular proporcionan un potente marco computacional para investigar estes procesos ao ofrecer unha descrición a nivel molecular e unha visión temporal das estruturas, dinámicas e interaccións dos péptidos.

Non obstante, as simulacións de Dinámica Molecular de sistemas que comprenden membranas lipídicas e múltiples péptidos cíclicos están severamente limitadas polo custo computacional, particularmente cando se estudan fenómenos colectivos como a autoensamblaxe, a inserción na membrana ou a agregación. Para superar estas limitacións, os enfoques de Dinámica Molecular \textit{Coarse-Grained} ofrecen unha alternativa práctica ao reducir o número de graos de liberdade e permitir simulacións en escalas de tempo máis longas e tamaños de sistema maiores. No caso dos \textit{D,L}-$\alpha$-péptidos cíclicos, o uso de modelos \textit{coarse-grained} introduce un desafío crítico: a súa autoensamblaxe réxese pola orientación do esqueleto e a direccionalidade das pontes de hidróxeno, ambas intrinsecamente ligadas á quiralidade. Os campos de forza \textit{coarse-grained} estándar, como Martini, non codifican explicitamente esta información, o que leva a xeometrías de amoreamento pouco realistas e descricións inexactas da ensamblaxe supramolecular. Ter en conta con precisión a quiralidade a nivel \textit{coarse-grained} tórnase, por tanto, esencial para modelar fielmente estes sistemas. Ao mesmo tempo, a adopción de representacións \textit{coarse-grained} permite estratexias de simulación sistemáticas e de alto rendemento, permitindo a exploración dun gran número de secuencias de péptidos cíclicos e escenarios de interacción coa membrana baixo condicións totalmente controladas e reproducibles. Tales estudos xeran conxuntos de datos grandes e homoxéneos que describen a estrutura, a dinámica e as interaccións coa membrana dos péptidos. A dispoñibilidade destes conxuntos de datos é particularmente relevante no contexto da Aprendizaxe Automática. Os modelos preditivos actuais para a actividade dos péptidos cíclicos dependen predominantemente de representacións baseadas na secuencia, mapeando directamente as secuencias de aminoácidos a resultados biolóxicos. Aínda que exitosos en moitos casos, estes enfoques a miúdo carecen de información explícita sobre os mecanismos moleculares subxacentes que gobernan as interaccións péptido-membrana. Pola contra, as simulacións de Dinámica Molecular proporcionan acceso a descritores ricos e fisicamente significativos---como enerxías de interacción, patróns de contacto, reordenamentos estruturais e propiedades dinámicas---que están directamente ligados á función.

Esta tese doutoral fundaméntase nesta sinerxía entre a simulación molecular e a intelixencia artificial. O obxectivo final é proporcionar novos métodos e ferramentas computacionais que apoien o deseño racional de péptidos cíclicos, particularmente para aplicacións antimicrobianas, mediante simulacións fisicamente informadas e predición baseada en datos. Para lograr isto, o traballo estrutúrase arredor dunha serie de obxectivos centrais:
\begin{enumerate}
    \item O desenvolvemento dun campo de forzas \textit{coarse-grained} que teña en conta a quiralidade e mellore a representación dos péptidos cíclicos á vez que permite simulacións a escala de microsegundos.
    \item A creación de ferramentas computacionais de código aberto para a xeración automatizada de estruturas tridimensionais e topoloxías de secuencias de péptidos cíclicos e nanotubos de péptidos cíclicos autoensamblados, co obxectivo de reducir as barreiras técnicas e democratizar o acceso ás simulacións de Dinámica Molecular destes sistemas.
    \item A aplicación de modelos \textit{coarse-grained} que teña en conta a quiralidade e ferramentas computacionais para realizar simulacións de Dinámica Molecular sistemáticas e de alto rendemento para investigar a autoensamblaxe de péptidos cíclicos e as súas interaccións con modelos de membrana.
    \item A extracción, análise e organización de descritores derivados de simulacións nunha base de datos de Dinámica Molecular pública e reutilizable, deseñada para apoiar a análise baseada en datos e o desenvolvemento de modelos.
    \item O desenvolvemento e aplicación de modelos de Aprendizaxe Automática adestrados tanto con datos experimentais como con descritores derivados da Dinámica Molecular para predicir a permeabilidade de membrana dos péptidos cíclicos e guiar estratexias de deseño informadas polo mecanismo.
\end{enumerate}

O primeiro gran logro desta tese foi o desenvolvemento de MA(R/S)TINI, unha nova adaptación do campo de forzas \textit{coarse-grained} Martini 2.2 que ten en conta a quiralidade, especificamente deseñado para modelar a autoensamblaxe de \textit{D,L}-$\alpha$-péptidos cíclicos en nanotubos de péptidos cíclicos autoensamblados. Este modelo supera varias limitacións fundamentais do campo de forzas orixinal: a sobreestimación das distancias interpeptídicas dentro dos nanotubos, a súa incapacidade para distinguir entre disposicións de lámina $\beta$ paralelas e antiparalelas, e a formación de estruturas autoensambladas con ángulos de rotación pouco realistas que resultarían en desaxustes doador-aceptor a nivel atomístico. Esta última limitación é particularmente crítica, xa que leva a unha sobreestimación artificial da ensamblaxe dos péptidos cíclicos. Para avaliar o rendemento de MA(R/S)TINI, seleccionáronse tres secuencias representativas de péptidos cíclicos: CP1 (\underline{R}R\underline{K}W\underline{L}W\underline{L}W, anfipático con actividade antimicrobiana probada), CP2 (\underline{L}L\underline{L}L\underline{L}L\underline{L}L totalmente hidrófobo), e CP3 (\underline{L}Q\underline{L}W\underline{L}W\underline{L}W, hidrófobo con capacidade para formar nanotubos transmembrana). Estes péptidos simuláronse en presenza de tres modelos de membrana diferentes: un de mamífero e dous de tipo bacteriano. Os resultados comparáronse sistematicamente cos obtidos para os mesmos sistemas utilizando o campo de forzas Martini 2.2 estándar.

As simulacións a gran escala realizadas a diferentes concentracións de péptidos cíclicos (20, 40, 60, 80 e 160 unidades) demostraron que MA(R/S)TINI producía patróns de autoensamblaxe similares aos observados con Martini 2.2, á vez que proporcionaba unha descrición máis fisicamente consistente das xeometrías de interacción. É importante destacar que o campo de forzas capturou as interaccións esperadas das diferentes secuencias de péptidos cíclicos con composicións de membrana específicas. Por exemplo, CP1 formou nanotubos preferentemente na fase acuosa cando se simulou co modelo de mamífero, e na superficie da bicapa lipídica cando se expuxo a un modelo bacteriano. A diferenza de Martini 2.2, MA(R/S)TINI diferenciou correctamente entre disposicións paralelas e antiparalelas e restrinxiu os ángulos de rotación a valores con sentido físico. As simulacións de metadinámica de dímeros de péptidos cíclicos confirmaron ademais que a introdución de sitios de interacción adicionais no esqueleto abordaba eficazmente as limitacións xeométricas do campo de forzas orixinal. Non obstante, os cálculos do Potencial de Forza Media mostraron unha enerxía de interacción case dúas veces maior en MA(R/S)TINI, o que probablemente estaba a levar a unha sobreestimación da autoensamblaxe.

O lanzamento dunha nova versión de Martini, Martini 3, que introduciu un conxunto redefinido de partículas de interacción, permitiu o refinamento deste enfoque mediante o desenvolvemento de MA(R/S)TINI 3. Esta nova versión deseñouse para conservar as vantaxes xeométricas do modelo que ten en conta a quiralidade mentres restaura as enerxías de interacción a valores máis próximos aos do campo de forzas orixinal. Para garantir unha comparación xusta entre MA(R/S)TINI 3 e o seu predecesor, empregáronse os mesmos protocolos de simulación e sistemas de proba. A Dinámica Molecular non sesgada demostrou que MA(R/S)TINI 3 preservaba a autoensamblaxe dependente da secuencia e melloraba as xeometrías de interacción, á vez que producía ensamblaxes supramoleculares máis dinámicas consistentes cunha redución na forza da interacción interpeptídica. As simulacións de metadinámica confirmaron a mellora respecto á versión orixinal ao reducir a esaxerada enerxía de interacción entre péptidos cíclicos para igualala á do campo de forzas orixinal. Ao preservar os puntos fortes de MA(R/S)TINI e corrixir as súas debilidades, MA(R/S)TINI 3 ofrece unha plataforma robusta para estudar nanotubos supramoleculares baseados en péptidos cíclicos e as súas interaccións coas membranas.

Co fin de facilitar o acceso ás simulacións de Dinámica Molecular de péptidos cíclicos e á visualización de estruturas tridimensionais, integráronse tanto os campos de forzas \textit{coarse-grained} orixinais como os adaptados, xunto con múltiples modelos atomísticos, nunha ferramenta como parte dun servidor web dispoñible publicamente baixo o nome CYCLOPEp. A ferramenta CYCLOPEpBuilder proporciona unha interface amigable para a creación de secuencias de péptidos cíclicos ou nanotubos de péptidos cíclicos autoensamblados. Esta ferramenta proporciona un control total sobre a estrutura, incluíndo: a selección da lonxitude da secuencia que vai de 4 a 12 aminoácidos, a personalización da orientación relativa entre pares de péptidos cíclicos, así como a rotación dunha secuencia relativa á súa veciña. Tras uns cálculos mínimos, a plataforma proporciona un conxunto de ficheiros que conteñen as estruturas tridimensionais do(s) péptido(s) cíclico(s) e do nanotubo ensamblado, xunto cos ficheiros de simulación necesarios para configurar unha simulación de Dinámica Molecular. Ao automatizar a configuración de estruturas complexas e minimizar as barreiras técnicas, CYCLOPEpBuilder democratiza o acceso aos métodos de simulación, promovendo a súa adopción dentro da investigación de péptidos cíclicos.

Baseándose nestas ferramentas, realizouse un estudo de DM a gran escala, resultando en CYCLOPEpDB, unha base de datos de simulacións de Dinámica Molecular \textit{coarse-grained} de sistemas péptido cíclico-membrana usando MA(R/S)TINI 3. A base de datos inclúe resultados para 3855 \textit{D,L}-$\alpha$-octapéptidos cíclicos únicos interactuando cun modelo de membrana bacteriana. Cada secuencia utilizouse para construír un nanotubo de oito membros que foi simulado en tres orientacións relativas á superficie da membrana. Unha configuración e un fluxo de traballo de simulación uniformes garantiron condicións consistentes a través dos experimentos. De cada traxectoria de 10 microsegundos, extraéronse 1006 descritores cuantitativos, cubrindo enerxías de interacción, número de contactos ou distancias, entre outros.

Como proba de concepto, seleccionouse un subconxunto de seis descriptores que capturan aspectos clave da interacción entre os péptidos cíclicos e a membrana a partir do conxunto completo de descriptores. A continuación, aplicouse a Análise de Compoñentes Principais para reducir a dimensionalidade dos datos a dúas compoñentes principais, que foron logo utilizadas para o agrupamento co algoritmo OPTICS. A análise da distribución dos seis descriptores iniciais a través dos tres clústeres identificados revelou distintas clases de interacción coa membrana, que foron categorizadas como Lipofílicas, Intermedias e Lipofóbicas. Unha análise posterior destas categorías revelou correlacións coa composición de aminoácidos e a carga neta: os péptidos cíclicos Lipofílicos estaban enriquecidos en residuos básicos pero tiñan preferentemente cargas netas próximas a cero, as secuencias Intermedias estaban maioritariamente moi cargadas positivamente cun perfil de aminoácidos equilibrado, mentres que os péptidos cíclicos Lipofóbicos eran ricos en residuos ácidos e portaban predominantemente unha carga neta negativa.

Para permitir as predicións de Aprendizaxe Automática baseadas nos datos de simulación, cada secuencia converteuse nun grafo, baseándose na súa topoloxía \textit{coarse-grained} e nos parámetros do campo de forzas. Os nodos do grafo codificaron características a nivel de partícula (masa, carga, coordenadas e valores de interacción de Lennard-Jones), mentres que as arestas representaron a conectividade dos enlaces e as lonxitudes dos mesmos. Unha Rede Neuronal de Grafos adestrada nesta representación das moleculas como grafos logrou un rendemento sobresaliente tanto na tarefa de clasificación cunha exactitude próxima ao 90\% como na tarefa de regresión predicindo unha puntuación de pseudopermeabilidade baseada na simulación cun \(R^2 \approx\)0.92 e un erro cadrático medio de \(0.51\). A arquitectura utilizada en combinación cos grafos baseados en MA(R/S)TINI 3 superou significativamente aos modelos de referencia, demostrando o poder de combinar simulacións de Dinámica Molecular coa Aprendizaxe Automática baseada en grafos. Ademais, todos os datos de simulación e os descritores extraídos puxéronse a disposición pública mediante a integración de CYCLOPEpDB na plataforma CYCLOPEp, permitindo a outros investigadores realizar as súas propias análises e adestrar modelos preditivos.

Paralelamente, o lanzamento de CycPeptMPDB, unha base de datos que contén datos experimentais de permeabilidade de membrana para péptidos cíclicos, permitiu o desenvolvemento de modelos de Aprendizaxe Automática adestrados directamente con medicións experimentais. O obxectivo deste estudo foi optimizar tanto o rendemento de clasificación como o de regresión aproveitando a intuición química na elección de representacións moleculares, estratexias de aumento de datos e arquitecturas de redes neuronais. En primeiro lugar, avaliáronse sistematicamente tres tipos de representacións de entrada:
\begin{enumerate}
    \item \textit{Sequential properties} que codifica a orde dos aminoácidos con características a nivel de residuo.
    \item Descritores fisicoquímicos a nivel de péptido cíclico.
    \item Codificación baseada en \textit{SMILES}.
\end{enumerate}
O modelo con mellor rendemento, que combinaba \textit{sequential properties} e características a nivel de péptido cíclico, foi adestrado con dúas estratexias diferentes de aumento de datos. O primeiro método explotou a circularidade das secuencias aplicando permutacións cíclicas, é dicir, desprazando a ordenación dos aminoácidos dentro da secuencia nunha posición, engadindo ao conxunto de datos todas as combinacións posibles que representaban cada secuencia. O segundo método tiña como obxectivo aumentar enormemente o número de secuencias mediante substitucións de aminoácidos. Dado que a substitución dun aminoácido por outro pode ter un grande impacto na actividade, especialmente en péptidos curtos, probáronse diferentes regras de substitución. Algúns conxuntos de datos foron deseñados para seren máis conservadores, buscando unha mellora no rendemento da predición, mentres que outros introduciron cambios disruptivos para confundir ao modelo e estudar o papel da mutación adecuada de aminoácidos. Aínda que as permutacións cíclicas mostraron unha lixeira mellora da exactitude preditiva, todos os conxuntos de datos de mutación estudados reduciron o rendemento do modelo orixinal. Observouse que as mutacións máis conservadoras produciron modelos máis exactos, salientando a importancia dunha selección coidadosa dos aminoácidos.

Finalmente, avaliaronse dúas modificacións arquitectónicas para mellorar a capacidade do modelo para capturar a natureza cíclica das secuencias. Nun primeiro enfoque, as capas convolucionais convencionais substituíronse por capas convolucionais cíclicas, especificamente deseñadas para ter en conta os patróns de peche de anel. Esta modificación produciu resultados moi similares aos do modelo orixinal, suxerindo que as capas convolucionais estándar poden capturar xa a información cíclica sen modificacións adicionais. Nun segundo enfoque, cada permutación cíclica dunha secuencia procesouse independentemente utilizando unha rede neuronal convolucional seguida dun módulo de Memoria a Curto e Longo Prazo, e os tensores de saída agregáronse para producir as predicións finais. A pesar dun lixeiro aumento no custo computacional, o modelo logrou rendementos preditivos comparables aos de redes neuronais profundas máis complexas publicadas na literatura, suxerindo que a calidade dos datos, máis que a complexidade do modelo, pode ser o factor limitante.

Para facilitar unha ampla accesibilidade e un uso práctico, o modelo experimental de permeabilidade con mellor rendemento despregouse como CYCLOPS, unha ferramenta web integrada na plataforma CYCLOPEp. CYCLOPS permite aos usuarios predicir a permeabilidade de membrana utilizando a secuencia do péptido cíclico como entrada. O construtor de secuencias soporta máis de 300 aminoácidos e permite a construción de péptidos cíclicos que van de 2 a 15 residuos con esquemas de ciclación tanto cabeza-cola como cabeza-cadea lateral.

En resumo, este traballo destaca a importancia das representacións moleculares quimicamente informadas e das estratexias de aumento de datos prudentes co fin de construír modelos exactos para a predición da permeabilidade de péptidos cíclicos. Ademais, proporciona unha plataforma de uso sinxelo que ofrece accesibilidade a ferramentas de Aprendizaxe Automática a investigadores sen experiencia.

En conclusión, esta tese doutoral proporciona CYCLOPEp, unha plataforma en liña e de código aberto que integra tres compoñentes complementarios destinados a acelerar a investigación en péptidos cíclicos. CYCLOPEpBuilder permite a xeración automatizada de estruturas tridimensionais e topoloxías de \textit{D,L}-$\alpha$-Péptidos Cíclicos e nanotubos de péptidos cíclicos autoensamblados, xa sexa para simulación ou visualización de estruturas. Están dispoñibles diferentes niveis de resolución, incluíndo os modelos \textit{coarse-grained} MA(R/S)TINI especificamente adaptados para reproducir as características dos péptidos cíclicos cun custo computacional reducido. CYCLOPEpDB proporciona unha extensa base de datos de simulacións de Dinámica Molecular \textit{coarse-grained} de \textit{D,L}-$\alpha$-octapéptidos cíclicos en presenza dun modelo de membrana bacteriana, ao nivel de resolución MA(R/S)TINI 3. A base de datos inclúe análises precalculadas e descritores derivados da simulación que poden ser utilizados directamente para a análise de datos e a Aprendizaxe Automática. Finalmente, CYCLOPS integra modelos de Aprendizaxe Automática adestrados con datos experimentais de permeabilidade de péptidos cíclicos para predicir a permeabilidade de membrana. O seu construtor de secuencias soporta máis de 300 monómeros, permitindo a construción de péptidos cíclicos de lonxitudes e esquemas de ciclación variables. Os resultados desta tese ofrecen unha plataforma sólida para o estudo de péptidos cíclicos utilizando ferramentas computacionais. Esperamos que CYCLOPEp sirva como un recurso valioso para a comunidade científica facilitando o acceso a ferramentas computacionais e datos, e apoiando a integración de técnicas de modelado computacional con experimentos de laboratorio cara ao deseño racional de péptidos cíclicos con actividade antimicrobiana.

































\subsection*{Resumen}
Una barrera de escasos nanómetros define la frontera física entre la vida y la materia inerte. A nivel celular, la vida depende por completo de la integridad, la dinámica y la funcionalidad de este límite conocido como la membrana celular. Se trata de una estructura compleja, flexible y dinámica que garantiza la estabilidad de la célula y regula el intercambio bidireccional de materia y energía con su entorno. Sin embargo, cuando se compromete la homeostasis de esta interfaz molecular, las consecuencias trascienden el nivel celular y desencadenan procesos patológicos severos.

Tradicionalmente, las patologías derivadas de estas alteraciones, de elevado impacto clínico y social (como el cáncer, las enfermedades infecciosas y las patologías degenerativas asociadas al envejecimiento), han sido estudiadas, diagnosticadas y tratadas como entidades clínicas completamente independientes. El punto de partida de esta tesis doctoral es el postulado de que estas tres grandes condiciones fisiopatológicas no están aisladas, sino que conforman los vértices de lo que se ha denominado el \textit{Triángulo CAIn} (Cáncer, Aging e Infección).

La revisión exhaustiva de la literatura pone de manifiesto la existencia de una estrecha interrelación bidireccional entre estos procesos. Así, se ha descrito ampliamente que las infecciones crónicas pueden inducir un estado inflamatorio persistente que favorece la aparición de transformaciones oncogénicas, como ocurre en el caso de \textit{Helicobacter pylori} y el cáncer gástrico. Del mismo modo, el envejecimiento se acompaña de un proceso de inmunosenescencia que incrementa de manera significativa tanto la incidencia de neoplasias como la susceptibilidad a enfermedades infecciosas. Más allá de estas asociaciones clínicas y fisiopatológicas, esta tesis plantea la existencia de un elemento común subyacente capaz de conectar estos tres ejes: la alteración anómala y sostenida de los perfiles lipídicos en dicha frontera nanométrica.

La membrana plasmática constituye una estructura compleja formada principalmente por lípidos, proteínas y carbohidratos conjugados, organizados en una arquitectura dinámica y altamente especializada. Entre sus componentes lipídicos, los fosfolípidos representan el fundamento estructural de la bicapa. En medios acuosos, estas moléculas anfifílicas se autoensamblan espontáneamente, un fenómeno gobernado por contribuciones entrópicas del solvente y estabilizado entálpicamente mediante interacciones intermoleculares entre las cadenas hidrofóbicas. En condiciones fisiológicas, la membrana plasmática de las células sanas de mamífero presenta una marcada asimetría transversal. La monocapa externa, expuesta al medio extracelular, está enriquecida en lípidos zwitteriónicos, principalmente fosfatidilcolina y esfingomielina, mientras que la monocapa interna contiene predominantemente fosfatidiletanolamina y fosfatidilserina, esta última de carácter aniónico. El mantenimiento de esta distribución asimétrica requiere un consumo energético continuo y resulta esencial para la viabilidad y funcionalidad celular.

No obstante, en los escenarios que conforman el triángulo CAIn, los mecanismos responsables del mantenimiento de esta asimetría colapsan. En las células tumorales se observa una tendencia hacia la pérdida de la organización transversal de la membrana, de modo que lípidos aniónicos normalmente confinados a la monocapa interna, como la fosfatidilserina, se externalizan hacia la superficie celular. Como consecuencia, la membrana de la célula cancerosa adquiere una carga neta más negativa. De manera análoga, una elevada densidad de carga negativa constituye también una característica distintiva de numerosos patógenos. Bacterias como \textit{Escherichia coli}, por ejemplo, presentan membranas enriquecidas en lípidos aniónicos, entre ellos el fosfatidilglicerol. Estas alteraciones composicionales y electrostáticas generan diferencias fisicoquímicas sustanciales entre las membranas de las células sanas y las de células tumorales o microorganismos patógenos. En consecuencia, la exposición de cargas negativas en la superficie celular emerge como una diana terapéutica de especial interés, ya que permite el desarrollo de estrategias capaces de discriminar selectivamente células enfermas frente a células sanas, minimizando así el daño sobre los tejidos no afectados.

Entre los mecanismos biológicos que explotan estas diferencias en la composición lipídica, destaca el empleado por el sistema inmunitario innato de multitud de organismos, incluidos los seres humanos. En esta primera línea de defensa participan los péptidos antimicrobianos (AMPs, por sus siglas en inglés). Estas moléculas son pequeñas cadenas de aminoácidos capaces de reconocer y destruir selectivamente a microorganismos (incluyendo bacterias Gram-positivas y Gram-negativas), hongos y células tumorales. Aunque presentan una gran diversidad estructural, los AMPs comparten propiedades fundamentales: suelen ser secuencias cortas, poseen una carga neta positiva (catiónica) y exhiben un marcado carácter anfifílico.

La actividad biológica de los AMPs está íntimamente ligada a su naturaleza fisicoquímica, la cual dirige su mecanismo de acción. En una fase inicial, las regiones catiónicas del péptido promueven una fuerte atracción electrostática hacia las membranas cargadas negativamente de las células objetivo. Una vez adsorbidos en la superficie celular, muchos de estos péptidos experimentan un cambio conformacional, adoptando frecuentemente una estructura helicoidal. Este plegamiento anfifílico amplifica la interacción con la bicapa: los residuos hidrofóbicos se insertan hacia el núcleo no polar de las colas lipídicas, mientras que los dominios polares se mantienen orientados hacia el exterior acuoso y las cabezas de los fosfolípidos.

Estas fuerzas combinadas provocan la acumulación progresiva de los AMPs en la membrana. Al alcanzar una concentración crítica, los péptidos se agregan y desencadenan efectos membranolíticos. Existen diversos mecanismos para explicar este proceso, como la formación de poros que permiten la fuga de material citoplasmático o la extracción directa de lípidos. En cualquier escenario, la disrupción de la integridad de la bicapa culmina inexorablemente en la muerte celular.

A diferencia de las terapias farmacológicas convencionales, la actividad de los AMPs no depende del reconocimiento de un receptor proteico específico, sino de las propiedades fisicoquímicas globales de la membrana diana. Este mecanismo de acción, basado en características estructurales generales más que en una interacción molecular específica, reduce significativamente la probabilidad de aparición de los mecanismos clásicos de resistencia antimicrobiana. En consecuencia, durante las últimas décadas se ha dedicado un considerable esfuerzo al desarrollo y optimización de estas moléculas. No obstante, su aplicación clínica continúa limitada por diversos obstáculos intrínsecos, entre los que destacan los elevados costes de producción, su susceptibilidad a la degradación y una eficacia frecuentemente reducida en condiciones \textit{in vivo}.


Superar estas limitaciones y optimizar el diseño racional de nuevos AMPs requiere una caracterización detallada de su interacción con las membranas celulares, prestando especial atención tanto a los cambios conformacionales inducidos tras la adsorción como a los mecanismos moleculares responsables de su actividad biológica. Tradicionalmente, el estudio de estos sistemas se ha abordado mediante estrategias experimentales multidisciplinares, con el objetivo de obtener una visión integrada del proceso y compensar las limitaciones inherentes a cada técnica. No obstante, incluso con los avances alcanzados en metodologías analíticas y estructurales, continúa siendo extremadamente difícil acceder simultáneamente a la evolución temporal completa del proceso de inserción y a una descripción conformacional con resolución atómica.
En este contexto, las aproximaciones computacionales se han consolidado como herramientas esenciales para complementar e interpretar los datos experimentales. En particular, la presente tesis doctoral aborda el estudio de los AMPs mediante una estrategia computacional multiescala, abarcando desde simulaciones de dinámica molecular clásica hasta metodologías emergentes basadas en computación cuántica.
La interacción entre AMPs y membranas lipídicas se ha estudiado tradicionalmente mediante simulaciones de Dinámica Molecular (MD, \textit{Molecular Dynamics}). Esta metodología permite describir la evolución temporal de un sistema molecular mediante la resolución numérica de las ecuaciones clásicas del movimiento. Aunque la MD no incorpora explícitamente los efectos cuánticos asociados al comportamiento electrónico, esta aproximación reduce de forma considerable el coste computacional, permitiendo el estudio de sistemas biomoleculares complejos en escalas espaciales y temporales inaccesibles para otros métodos atomísticos. Como consecuencia, la dinámica molecular se ha convertido en una herramienta fundamental para investigar la organización, estabilidad y dinámica de membranas biológicas y sus interacciones con biomoléculas.

Pese a que la complejidad de los sistemas simulados mediante MD ha crecido en paralelo a la capacidad de cálculo de los ordenadores modernos, siguen existiendo limitaciones. El elevado coste computacional y las altas barreras de energía restringen los fenómenos de interacción péptido-membrana que pueden ser capturados espontáneamente. Para superar este obstáculo, esta tesis recurre a técnicas de muestreo avanzado, como la Metadinámica. Estos métodos permiten forzar la exploración del sistema para reconstruir con precisión los paisajes de energía libre y cuantificar variables críticas, como la orientación espacial y el momento dipolar hidrofóbico del péptido al penetrar en diferentes modelos de membrana simplificados (sana, bacteriana y tumoral).

Más allá de la caracterización biofísica, la creación de nuevas secuencias terapéuticas exige un paso adicional. En la actualidad, la inmensa cantidad de datos generados en la investigación se explota frecuentemente para entrenar modelos predictivos estadísticos con el fin de desarrollar nuevas terapias. Sin embargo, en el diseño de AMPs, los enfoques puramente predictivos a menudo carecen de interpretabilidad física y presentan serias dificultades para evaluar secuencias en interfaces heterogéneas, como la frontera agua-lípido.

Por ello, alejándose de los modelos predictivos de caja negra, esta tesis postula un marco de estudio y diseño basado estrictamente en las interacciones físicas del sistema. Ante la necesidad de acelerar la exploración del vasto espacio de secuencias, un objetivo inabarcable mediante Dinámica Molecular debido a su elevado coste computacional, se propone el uso de algoritmos de optimización heurística y de tecnologías emergentes como computación cuántica (como VQE y QAOA).

Esta transición algorítmica no solo permite evaluar los péptidos de forma rápida mediante Hamiltonianos directamente condicionados por el entorno interfacial, sino que culmina en la reformulación del reto hacia la optimización combinatoria de secuencias. Este paso adapta el problema a la naturaleza inherentemente discreta del hardware cuántico, maximizando su rendimiento y viabilidad.

Por un lado, esta tesis emplea simulaciones de Dinámica Molecular y técnicas de metadinámica con el objetivo de alcanzar una comprensión atomística profunda de los mecanismos responsables de la selectividad interfacial. Por otro, frente al desafío asociado al descubrimiento de nuevos candidatos terapéuticos, desarrolla un marco algorítmico rápido, estructurado y potencialmente escalable para la exploración racional de interacciones membrana-péptido. De este modo, el trabajo trasciende la mera caracterización estructural y aporta herramientas orientadas al diseño racional de estrategias terapéuticas dirigidas frente a membranas patológicas.
Asimismo, la incorporación de metodologías inspiradas en computación cuántica establece una base conceptual y algorítmica preparada para futuros escenarios de ventaja cuántica práctica. En la medida en que la evolución del hardware cuántico permita abordar sistemas de mayor complejidad con costes computacionales competitivos, este tipo de aproximaciones podría ofrecer ventajas significativas frente a los métodos clásicos en problemas de optimización y exploración del espacio conformacional biomolecular.

El capítulo de Introducción profundiza en la estructura de las membranas celulares y su relevancia fisiopatológica. Continúa con una exposición sobre los péptidos antimicrobianos (AMPs), con especial énfasis en aquellos dirigidos a membrana, y las técnicas experimentales empleadas en su estudio. Posteriormente, se explica cómo la MD y las técnicas de muestreo avanzado, como la Metadinámica, complementan el estudio experimental abordando detalles atómicos y termodinámicos que las metodologías in vitro no pueden capturar de forma continua. El capítulo finaliza exponiendo las vías computacionales actuales para el desarrollo de nuevos AMPs, planteando la integración de estos conocimientos físicos con los recientes avances en computación clásica y cuántica.

Esta introducción sirve como base teórica para formular los Objetivos de este trabajo. De forma global, la tesis persigue enlazar la caracterización biofísica profunda de la interacción péptido-membrana con algoritmos de optimización cuántica y clásica, trazando una evolución metódica: desde el desafío de predecir conformaciones espaciales hasta la formulación de un diseño racional de secuencias. Para alcanzar esta meta, se establecen los siguientes objetivos específicos:
\begin{itemize}
    \item Primer objetivo: Analizar la interacción de distintos AMPs con diversos modelos de bicapas lipídicas (sanas y patológicas) mediante simulaciones de MD y Metadinámica. El fin es identificar y cuantificar descriptores físicos rigurosos, tales como la orientación del momento dipolar hidrofóbico transversal o las energías libres de adsorción, que dictan el mecanismo de acción y la selectividad de estas moléculas.
    \item Segundo objetivo: Evaluar la viabilidad de la computación cuántica para predecir la conformación de péptidos en interfaces lipídicas. Para ello, se implementan algoritmos híbridos, como el Variational Quantum Eigensolver (VQE), empleando modelos de red (\textit{lattices}) para aproximar el espacio espacial. Este paso aborda el reto fundamental de mapear un problema de naturaleza inherentemente continua (el plegamiento tridimensional) sobre el hardware intrínsecamente discreto de los ordenadores cuánticos actuales.
    \item Tercer objetivo: Pivotar hacia un enfoque de optimización de secuencias, un problema de naturaleza puramente discreta (la selección combinatoria de aminoácidos), para maximizar la idoneidad y el rendimiento del hardware cuántico. Esto se materializa en el desarrollo de un Hamiltoniano de diseño condicionado por el entorno interfacial. Su resolución mediante algoritmos como QAOA y heurística clásica permite generar secuencias terapéuticas \textit{de novo} orientadas a dianas específicas. Como resultado final, este marco físico se consolida en una plataforma de código abierto (Helix Design Studio), democratizando el diseño de péptidos en entornos heterogéneos.
\end{itemize}

Para alcanzar los objetivos planteados, esta tesis utiliza un enfoque computacional multiescala. Las simulaciones de MD atomística detallan las interacciones moleculares, pero su coste computacional limita el estudio de fenómenos colectivos como el autoensamblaje o la agregación peptídica. Para abordar escalas espaciales y temporales mayores, se emplean modelos de Dinámica Molecular de grano grueso (\textit{Coarse-Grained}, CG). Aunque esta reducción de resolución conlleva la pérdida de interacciones atómicas explícitas y obliga a mantener fija la estructura secundaria, la disminución de grados de libertad proporciona la velocidad computacional necesaria para, en combinación con técnicas de muestreo mejorado como la metadinámica, lograr una exploración extensa y eficiente del espacio conformacional global del sistema.

A nivel estructural, el documento avanza desde la caracterización biofísica hasta el diseño computacional. El capítulo de Metodología detalla los procedimientos empleados, comenzando por los fundamentos teóricos de la MD y las técnicas de muestreo avanzado, como la Metadinámica. Seguidamente, se describen los protocolos de simulación, la composición de los distintos modelos de membrana y los parámetros aplicados en cada etapa.

Posteriormente, se detallan las técnicas de análisis utilizadas sobre las trayectorias de simulación. Estas incluyen la caracterización de la orientación espacial de los péptidos en la interfaz, mediante el análisis de los ángulos de inclinación (\textit{tilt}) y rotación azimutal (\textit{spin}), y la cuantificación de los contactos específicos residuo-lípido. Asimismo, se describe el núcleo central del estudio termodinámico basado en el uso de métodos de muestreo avanzado: la reconstrucción de los perfiles de energía libre. Estos perfiles permitieron calcular las variaciones de energía libre ($\Delta$ G) asociadas a la inserción interfacial, cuantificando así los paisajes energéticos que gobiernan la selectividad de los péptidos.

El capítulo concluye detallando los algoritmos de optimización, tanto clásicos como cuánticos, formulados para resolver los Hamiltonianos de diseño en medios no homogéneos, como es la compleja interfaz agua-lípido. La integración de estos modelos computacionales permite la generación de novo de secuencias peptídicas optimizadas para adoptar estructuras \textit{$\alpha$-helicoidales} en entornos de membrana, proporcionando una solución robusta al reto de la especificidad terapéutica.


El primer bloque de resultados de esta tesis se deriva de un estudio bibliográfico sistemático de la literatura científica actual, centrado en la intersección de la inmunología, la oncología y la gerontología. El principal hito de esta fase fue la conceptualización y propuesta del \textit{Triángulo CAIn} (Cáncer, Aging e Infección), un marco teórico que unifica estas tres condiciones bajo un nexo fisiopatológico común.

A partir del análisis de los datos reportados en la literatura, se identificó que el estado inflamatorio crónico y persistente es el motor que vincula de forma bidireccional los tres vértices del triángulo. El estudio permitió concluir que:

\begin{itemize}
\item Correlación clínica: Existe una relación mecánica probada donde la inflamación derivada de infecciones crónicas actúa como precursor de procesos oncogénicos, mientras que la inmunosenescencia propia del envejecimiento facilita tanto la susceptibilidad a patógenos como la progresión tumoral.

\item El lipidoma como denominador común: El análisis bibliográfico reveló que, independientemente del origen de la patología, existe una alteración sistemática y convergente en la composición lipídica de las membranas celulares afectadas. El hallazgo más significativo fue la identificación de la pérdida de asimetría de la bicapa como un marcador universal.

\item Diana teranóstica: Se determinó que la exposición de lípidos aniónicos, principalmente fosfatidilserina, en la cara externa de la membrana plasmática no solo es una consecuencia del daño celular, sino que constituye una diana física específica para el desarrollo de nuevas estrategias teranósticas.
\end{itemize}

Este resultado bibliográfico permite concluir que la membrana plasmática deja de ser un mero límite celular para convertirse en una plataforma de información sobre el estado patológico del organismo. Esta sistematización teórica justifica el desarrollo de los capítulos posteriores de esta tesis, centrados en el diseño de AMPs capaces de reconocer de forma selectiva estas firmas lipídicas identificadas.


El segundo bloque de resultados expone la caracterización biofísica cuantitativa de la interacción entre AMPs y diferentes modelos de membrana. Para superar las limitaciones de la simulación clásica, se emplearon simulaciones de MD acopladas a técnicas de muestreo avanzado (Metadinámica). Esta aproximación permitió reconstruir las superficies de energía libre (FES) y calcular el potencial de fuerza media (PMF) asociados al proceso de adsorción e inserción peptídica.

Se analizaron comparativamente modelos de membranas sanas (zwitteriónicas, representativas de mamífero) frente a membranas patológicas (aniónicas, representativas de bacterias y células tumorales). Los cálculos termodinámicos demostraron que la selectividad de los AMPs está directamente gobernada por los perfiles de composición lipídica. Específicamente, se cuantificó que las membranas aniónicas reducen drásticamente las barreras energéticas para la inserción del péptido. Los resultados mostraron que, en estos entornos, la interacción electrostática inicial guía al péptido hacia una orientación interfacial óptima, permitiendo que el componente transversal del momento dipolar hidrofóbico penetre en el núcleo lipídico de forma termodinámicamente favorable.

Por el contrario, los perfiles de energía libre obtenidos sobre membranas de mamífero sanas revelaron la existencia de altas barreras energéticas que penalizan severamente la inserción profunda del péptido. En estos sistemas, los datos confirmaron que el péptido se mantiene confinado predominantemente en un estado de adsorción superficial transitoria. Mediante este análisis, se logró extraer y cuantificar los descriptores físicos rigurosos (FES, PMF y orientación espacial) que rigen la interacción péptido-membrana en función de la carga superficial objetivo.


El tercer bloque de resultados detalla la aplicación de algoritmos de computación cuántica para predecir el plegamiento de péptidos antimicrobianos en la interfaz agua-lípido. Para abordar este problema, se discretizó el espacio conformacional continuo mediante modelos de red (\textit{lattices} tridimensionales), formulando un Hamiltoniano físico que codifica explícitamente la hidrofobicidad de los residuos, las interacciones electrostáticas y el efecto limitante del entorno interfacial de la membrana.

Se evaluó el rendimiento del algoritmo híbrido VQE mediante su ejecución en simuladores cuánticos. Los resultados obtenidos confirmaron que el modelo VQE es capaz de identificar el estado fundamental del sistema. Las simulaciones lograron reproducir con exactitud las conformaciones plegadas de mínima energía (estructuras anfipáticas óptimas) inducidas por la presencia de la bicapa, obteniendo una concordancia directa con los métodos de resolución clásica exacta.

Adicionalmente, el estudio cuantificó el coste computacional y los requisitos de codificación sobre el hardware cuántico actual (dispositivos NISQ). Los datos métricos extraídos demostraron que el mapeo de trayectorias espaciales sobre redes discretas genera un incremento drástico en el número de qubits necesarios a medida que aumenta la longitud de la cadena peptídica. Estos resultados permitieron establecer empíricamente los límites actuales de escalabilidad del enfoque geométrico para la predicción de conformaciones estructurales complejas.


El cuarto y último bloque de resultados presenta el desarrollo de un marco físico interpretable para el diseño \textit{de novo} de secuencias peptídicas $\alpha$-helicoidales condicionadas por su entorno. Como respuesta a las limitaciones de escalabilidad de los modelos de predicción conformacional en red, el problema geométrico continuo se reformuló hacia la optimización discreta de secuencias (selección combinatoria de aminoácidos), un enfoque inherentemente más idóneo para su resolución algorítmica debido a la alta discretización que necesitan los actuales dispositivos cuánticos.

Para llevar a cabo este diseño, se desarrolló un Hamiltoniano basado estrictamente en principios físicos. Este modelo integra en una única función de energía las variables termodinámicas críticas identificadas en las etapas anteriores: la propensión helicoidal intrínseca, las interacciones inter-residuo, la electrostática, la exposición anisotrópica al disolvente (momento dipolar hidrofóbico) y la geometría específica de la interfaz lipídica.

Con el objetivo de permitir la comparación directa entre secuencias de diferente longitud y frente a distintos entornos de membrana, el modelo implementó una normalización estadística rigurosa. Todas las contribuciones energéticas fueron reescaladas y expresadas como Z-scores respecto a conjuntos de secuencias aleatorias de referencia, estableciendo un paisaje de diseño normalizado e independiente del tamaño del sistema.

La resolución de este Hamiltoniano, validada mediante métodos de heurística clásica y algoritmos de optimización cuántica (QAOA), demostró la viabilidad de generar secuencias terapéuticas optimizadas específicamente para dianas interfaciales patológicas. Finalmente, la integración de estos modelos físicos y algoritmos de evaluación cristalizó en el desarrollo del Helix Design Studio. Esta plataforma de software de código abierto consolida la base teórica y computacional de la tesis, proporcionando una herramienta accesible y automatizada para el diseño racional de péptidos en entornos heterogéneos.



El cuerpo de este manuscrito concluye con un capítulo dedicado a las principales conclusiones derivadas de las distintas etapas de la investigación. En su conjunto, esta tesis ha abordado el estudio de la selectividad y el diseño de sistemas péptido-membrana desde un enfoque multiescala, fundamentado estrictamente en principios físicos, que enlaza la biofísica computacional con algoritmos de optimización avanzados.

En una primera fase, la revisión sistemática de la literatura permitió establecer el marco teórico del \textit{Triángulo CAIn} (Cáncer, Envejecimiento e Infección). Este estudio identificó la pérdida de asimetría lipídica y la consiguiente exposición de cargas negativas en la superficie celular como un nexo fisiopatológico común y una diana teranóstica de carácter universal.

Para comprender el reconocimiento de esta diana a nivel atómico, se emplearon simulaciones de Dinámica Molecular y Metadinámica. La aplicación de estas técnicas permitió cuantificar los descriptores termodinámicos, como las superficies de energía libre (FES) y el potencial de fuerza media (PMF), que gobiernan la selectividad de los péptidos antimicrobianos (AMPs). De este modo, se demostró a nivel físico cómo las membranas aniónicas reducen drásticamente las barreras energéticas para la inserción del momento dipolar hidrofóbico en el núcleo lipídico.

A partir de la caracterización de estos descriptores, se evaluó la viabilidad de la computación cuántica para predecir el plegamiento interfacial de los péptidos. Si bien algoritmos híbridos como el VQE demostraron su eficacia para identificar conformaciones óptimas en modelos discretos de red, los resultados también evidenciaron las actuales limitaciones de escalabilidad del hardware cuántico (NISQ) al intentar mapear un problema espacial de naturaleza continua.

En respuesta a esta barrera técnica, el enfoque metodológico pivotó hacia un escenario computacionalmente más idóneo y de naturaleza inherentemente discreta: la optimización combinatoria de secuencias. Para ello, se formuló un Hamiltoniano de diseño condicionado por el entorno que integra variables físicas rigurosas (electrostática, propensión helicoidal e hidrofobicidad) y cuya resolución es viable tanto mediante heurística clásica como a través de algoritmos cuánticos (QAOA).

Finalmente, todo este marco teórico y matemático ha cristalizado en el desarrollo del Helix Design Studio. Esta plataforma de software de código abierto automatiza la generación \textit{de novo} de secuencias peptídicas, evaluadas y normalizadas mediante Z-scores. En conclusión, los resultados de esta tesis demuestran que la integración de una descripción termodinámica precisa con las nuevas arquitecturas de computación permite superar la dependencia de los modelos estadísticos de \textit{caja negra}, estableciendo un marco metodológico sólido, físico y escalable para el diseño racional de nuevas terapias dirigidas a membrana.










